\section{Related Work}
\label{sec:related}

\para{Chain-of-thought prompting.} \citet{wei2022chain} introduced \chainofthought prompting, demonstrating that including intermediate reasoning steps in few-shot exemplars substantially improves LLM performance on arithmetic, commonsense, and symbolic reasoning tasks. \citet{kojima2022large} showed that the simple prompt ``Let's think step by step'' elicits useful reasoning even without exemplars (zero-shot CoT). \citet{wang2023selfconsistency} improved upon greedy CoT decoding by sampling multiple reasoning paths and selecting the most consistent answer via majority voting. These foundational works establish step-by-step decomposition as the dominant CoT format, but do not explore alternative reasoning structures.

\para{Structural variants of CoT.} Several works have explored non-linear reasoning structures. Tree of Thoughts \citep{yao2023tree} extends CoT from linear chains to tree structures, enabling exploration and backtracking. Chain of Code \citep{li2023chain} uses executable code as the reasoning medium. Buffer of Thoughts \citep{yang2024buffer} stores and retrieves high-level reasoning templates. Contrastive CoT \citep{chia2023contrastive} augments demonstrations with both correct and incorrect reasoning chains. These approaches modify the \emph{structure} of reasoning (linear vs.\ branching, natural language vs.\ code) but maintain step-by-step logical decomposition within each path.

\para{Narrative reasoning in LLMs.} Most directly relevant to our work, \citet{zhang2024narrative} proposed Narrative-of-Thought (NOT), which prompts LLMs to generate temporally grounded narratives before producing temporal graphs. NOT achieved 16--71\% F1 improvements over standard prompting on temporal graph generation benchmarks. However, NOT was evaluated only on temporal reasoning tasks with structured graph outputs---a domain where narrative structure has a natural advantage. Our work tests whether the benefits of narrative-form reasoning generalize to standard QA benchmarks across diverse reasoning domains.

\para{Why does CoT work?} A growing body of work examines the mechanisms behind CoT's effectiveness. \citet{schaeffer2023invalid} showed that even \emph{logically invalid} CoT prompts achieve gains comparable to valid ones, suggesting that CoT may work through computational scaffolding rather than genuine logical reasoning. \citet{arcuschin2025chain} demonstrated that CoT reasoning is not always faithful---the stated reasoning may not reflect the model's actual computation. \citet{chen2025towards} surveyed long CoT reasoning and identified deep reasoning, exploration, and reflection as key characteristics. Our finding that narrative-form and step-by-step CoT achieve equivalent performance adds to this literature by showing that the \emph{format} of intermediate reasoning is not a critical factor for state-of-the-art models.

\para{Implicit and internalized reasoning.} \citet{deng2023implicit} distilled explicit CoT reasoning into a model's internal representations, enabling reasoning without explicit intermediate tokens. Diffusion of Thoughts \citep{ye2024diffusion} generates all reasoning tokens simultaneously rather than sequentially. These works suggest that the surface form of reasoning may be less important than the underlying computation it enables---a view consistent with our finding of format equivalence between narrative and step-by-step CoT.
